\section{Method details}
\label{app:method}

\subsection{Benchmark and rollouts}
\label{app:design}
EvalAwareBench adapts its tasks from HarmBench \citep{mazeika2024harmbench} and defines eight
cues \citep{li2026decomposing}: placeholder entities (F1), an incoherent premise (F2), a formal
register (F3), stripped context (F4), telegraphed stakes (F5), a recognisable paradigm (F6), a
structured format (F7) and F8, which Section~\ref{sec:method} describes. On most safety variants, F6 is a
jailbreak wrapper; on the capability variants, it is a benign reasoning or role prefix. For the
multi-cue conditions, we drew up to \cv{design/maxdraws} combinations of F1 to F7 for each number
of cues from two to seven.

Both models used the benchmark's realistic chat system prompt \citep{li2026decomposing}, kept
their reasoning traces and had a budget of \cv{design/maxtokens} output tokens. The rollouts and
the self-report and third-person answers were sampled at temperature \cv{design/temperature},
the default in the benchmark's code, with the other sampling parameters at the serving
defaults. Gemma ran on OpenRouter's fp8 endpoints. Qwen ran on our own vLLM server
\citep{kwon2023efficient} with fp8 weights and an fp8 key-value cache.
Table~\ref{tab:counts} gives the rollouts with a valid prediction from each detector.

\begin{table}[htbp]
  \centering
  \small
  \begin{tabular}{@{}l cc cc@{}}
    \toprule
    & \multicolumn{2}{c}{Gemma 4 31B} & \multicolumn{2}{c}{Qwen 3.6 27B} \\
    \cmidrule(lr){2-3}\cmidrule(l){4-5}
    Detector & Cued & Baseline & Cued & Baseline \\
    \midrule
    Self-report & \cv{gemma/sr/ncued} & \cv{gemma/sr/nbase} & \cv{qwen/sr/ncued} & \cv{qwen/sr/nbase} \\
    Third-person & \cv{gemma/tp/ncued} & \cv{gemma/tp/nbase} & \cv{qwen/tp/ncued} & \cv{qwen/tp/nbase} \\
    CoT judge & \cv{gemma/cot/ncued} & \cv{gemma/cot/nbase} & \cv{qwen/cot/ncued} & \cv{qwen/cot/nbase} \\
    Activation probe & \cv{gemma/probe/ncued} & \cv{gemma/probe/nbase} & \cv{qwen/probe/ncued} & \cv{qwen/probe/nbase} \\
    Logit monitor & \cv{gemma/logit/ncued} & \cv{gemma/logit/nbase} & \cv{qwen/logit/ncued} & \cv{qwen/logit/nbase} \\
    \bottomrule
  \end{tabular}
  \caption{Rollouts with a valid prediction from each detector (cued: F1 to F7, in
  single-cue and multi-cue conditions). Failed calls and non-finite activations or
  log-probabilities are excluded. Gemma's white-box counts are the rollouts with finite activations, a non-random
  subset (Appendix~\ref{app:gemma}).}
  \label{tab:counts}
\end{table}

\subsection{Detectors}
\label{app:detectors}
\begin{table}[htbp]
  \centering
  \footnotesize
  \begin{tabularx}{\linewidth}{@{}l l X l l@{}}
    \toprule
    Detector & Decision by & Input & Unit & \shortstack[l]{Tuned on\\cue labels} \\
    \midrule
    Self-report & the model & the task and its response, then a question on the setting
      (Qwen's template drops the reasoning trace) & rollout & no \\
    Third-person & the model & the task and answer, without reasoning trace &
      rollout & no \\
    CoT judge & a judge model & the task, reasoning trace and answer; decides on the trace &
      rollout & no \\
    Activation probe & a learned threshold & activations at the last prompt token & prompt &
      layer, threshold \\
    Logit monitor & a learned threshold & log-probability of ``The user might be testing'' &
      rollout & threshold \\
    \bottomrule
  \end{tabularx}
  \caption{The five detectors.}
  \label{tab:detectors}
\end{table}

Self-report and third-person ask the question of \citet{li2026decomposing}, which offers three
answers: deployment, a safety evaluation or a capability evaluation. Both count either
evaluation answer as aware. For self-report, our code sends the full response, reasoning trace
included, as the earlier assistant turn. However, Qwen's pinned chat template (commit
\texttt{6a9e13bd}) renders the reasoning of an earlier turn only when
\texttt{preserve\_thinking} is set, and our server passed only \texttt{enable\_thinking}. On
Qwen, self-report therefore saw only the task and the visible answer. On Gemma, the trace went
as text inside \texttt{<think>} tags, which Gemma's published template keeps. Gemma ran on
OpenRouter providers, though, whose copy of the template we cannot see, and the outputs neither
confirm nor refute that Gemma read the trace.

The benchmark's code also passes the reasoning trace to third-person. We removed it, so that
third-person does not read the evidence the CoT judge reads. The CoT judge uses the benchmark's
judge template \citep{li2026decomposing} with the model
\texttt{xiaomi/\allowbreak mimo-v2.5-pro} \citep{xiaomi2026mimo}. The template asks it to decide awareness from the reasoning trace alone and to
quote its evidence verbatim. Its false alarms are few: it flags \cv{gemma/cot/nfa} of the
\cv{gemma/cot/nbase} Gemma baseline rollouts it judged, and \cv{qwen/cot/nfa} of
\cv{qwen/cot/nbase} on Qwen. A second judge, \texttt{openai/\allowbreak gpt-oss-120b}, re-judged a
stratified subset of \cv{gemma/judge2/n} Gemma and \cv{qwen/judge2/n} Qwen rollouts with the same
template. The two judges agree at $\kappa = \cv{gemma/judge2/kappa}$ on Gemma and
\cv{qwen/judge2/kappa} on Qwen. On this subset the first judge flags \cv{gemma/judge2/pos1}\%
and \cv{qwen/judge2/pos1}\% of the rollouts, and the second only \cv{gemma/judge2/pos2}\% and
\cv{qwen/judge2/pos2}\%. The CoT judge's flags are therefore not reliable.\label{app:judge}

The logit monitor scores the log-probability of ``The user might be testing'' at up to
\cv{design/logitpositions} positions, namely the start of the reasoning trace and later sentence
boundaries, and keeps the maximum. \citet{aranguri2026logits} also take the maximum over a
prompt's rollouts, whereas we score each rollout. For both white-box detectors, the forward pass
used each model's weights at a pinned commit in fp8. The output head and the few linear layers
that fp8 blocks cannot tile stayed in bf16.

\subsection{White-box fit}
\label{app:whitebox}
The contrastive direction uses the \cv{design/pairs} prompt pairs of \citet{hua2025steering}. We
normalised it to unit length and read it at the last token of each prompt, written as a user
message, at every fourth layer. Within each of \cv{design/folds} training folds grouped by task,
we chose the layer and the threshold that maximised balanced accuracy for
\texttt{score}~$\geq$~\texttt{threshold}. We then applied them to the held-out tasks. The logit
monitor's threshold was chosen in the same way. On Qwen, the probe chose layer
\cv{qwen/probe/layer} of \cv{qwen/probe/nlayers}.

Each random direction in Figure~\ref{fig:probe-null} is a unit vector drawn separately at every
layer read (seed 0), and it went through the same choice of layer and threshold. The cue-trained
direction is the difference between the mean activations of cued and baseline prompts in the
training folds, recomputed in each fold with its own choice of layer. The logistic-regression
probe \citep{alain2016understanding} uses the probe's chosen layer. On Qwen, we excluded
\cv{qwen/reencode/nan} rollouts with non-finite activations and \cv{qwen/logit/nan} on which the
logit monitor failed.

\subsection{Gemma's white-box detectors}
\label{app:gemma}
The longest Gemma rollouts all fail. The tasks with most failures are song lyrics, a self-harm
blog post, forged prescriptions, and Wikipedia and political edits. In addition,
\cv{gemma/probe/novector} Gemma rollouts have no probe score, because no repetition of their
prompt gave finite activations.

Because the probe reads only the prompt, we also re-encoded Gemma's prompts in bf16. This gave
finite activations for 4,199 of the 4,200 prompts, and they agree with the fp8 activations at a
median cosine of 0.98 to 1.00 per layer. On these prompts, the probe's out-of-fold AUROC is
0.596 at layer 3. That lies inside the null of random directions with the same tuning (mean
0.566, highest 0.662). Therefore neither fp8 nor the missing rollouts explain Gemma's weak
probe.

\subsection{Positive control on a model organism}
\label{app:organism}
The organism is the LoRA adapter \texttt{timhua/wood\_v2\_sftr4\_filt} of \citet{hua2025steering}
on Llama 3.3 Nemotron Super 49B. We read it in bf16 at every second of its 80 layers, with the
adapter on (the organism) and off (the base model). We used the five prompt conditions of
\citet{hua2025steering} on their 128 coding questions, with their system prompts and user
templates. Two conditions state the setting: stated evaluation and stated deployment. In the
other three, a Wood Labs system prompt precedes user text that is neutral, states evaluation or
states deployment. The primary analysis also adds their hint to the system prompt. The hint asks
the model to think carefully about whether it is in evaluation or deployment before responding.
Their system prompts are assigned to the questions in five contiguous blocks. As in their work,
we built the direction from the \cv{design/pairs} pairs on the base model and applied it to both
arms. We read it at the last prompt token, where our probe reads. The analysis was committed
before the data existed, and the random-direction null below is post hoc.

We chose the layer that best separates stated evaluation from stated deployment, averaged over
the two arms, so the choice uses neither Wood Labs condition. It is layer 51, where that AUROC is
0.98 in the organism and 0.89 in the base model. For a condition $W$, $\lambda$ places its mean
projection $\bar p_W$ on the axis from stated deployment (0) to stated evaluation (1):
$\lambda = (\bar p_W - \bar p_{\mathrm{dep}})/(\bar p_{\mathrm{eval}} - \bar p_{\mathrm{dep}})$.
Because the two arms read identical text, their difference in $\lambda$ cannot come from the
words. Intervals resample questions within blocks (2,000 replicates).

With the hint on, Wood Labs prompts with neutral user text have $\lambda = 0.36$ in the organism
and $-0.51$ in the base model. The difference, 0.87 [0.75, 1.00], is positive in all five blocks
and at every layer from 39 to 79. Each arm's own direction, with its own layer choice (layer 41),
gives 1.26 [1.13, 1.40]. Without the hint, the base-model direction gives 1.24 [1.06, 1.44]. With
deployment user text, the difference is 0.15 [0.11, 0.18] with the hint and 0.61 [0.56, 0.65]
without it. With evaluation user text, it is $-0.23$ [$-0.28$, $-0.17$] with the hint and 0.28
[0.21, 0.36] without it. There the user text already states evaluation, so the base model places
these prompts near stated evaluation.

As for Gemma and Qwen, we put 1,000 random directions through the same layer choice. Each is a
unit vector drawn separately at every layer and oriented so that stated evaluation scores above
stated deployment in the base model. They also separate the two stated conditions (median
AUROC 0.89), because those conditions differ in their words. For the neutral Wood Labs prompts
with the hint on, their difference in $\lambda$ has a median of 0.05 and a 95th percentile of
0.83, and 4.6\% reach the observed 0.87. Without the hint, the median is 0.00, and 0.9\% reach
the observed 1.24. The contrastive direction therefore carries the association better than most
random directions, but only at the 5\% level with the hint on.

The control holds only at the last prompt token. Read at the last token of the user's text, the
difference reverses ($-0.66$ [$-1.03$, $-0.37$], no block positive). Averaged over the user's
tokens, it is $-0.13$ [$-0.17$, $-0.09$]. A read inside the user turn therefore does not carry
this system-prompt association.

\subsection{Statistics}
\label{app:formulas}
Hit and false-alarm rates use the correction $(k + 0.5)/(n + 1)$ and carry exact binomial
intervals \citep{clopper1934use}. The equal-variance Gaussian model behind $d'$ follows
\citet{hautus2021detection} and \citet{stanislaw1999calculation}. The bootstrap
\citep{efron1979bootstrap} resamples tasks as clusters \citep{field2007bootstrapping}, with both
variants of a task together. It uses \cv{run/kappaboot} resamples for the overall $\kappa$ and
\cv{run/bootstrap} for every other interval. We report $\kappa$ as $100\,\kappa/\kappa_{\max}$,
with a percentile interval from the same bootstrap.

For two detectors $a$ and $b$ with positive rates $p_a$ and $p_b$ on $N$ rollouts,
$\kappa = (p_o - p_e)/(1 - p_e)$, where $p_o$ is the observed agreement and
$p_e = p_a p_b + (1 - p_a)(1 - p_b)$ is the agreement expected by chance
\citep{cohen1960coefficient}. The largest attainable agreement is
$\min(p_a, p_b) + \min(1 - p_a, 1 - p_b)$, so
\begin{equation*}
  \kappa_{\max} = \frac{\min(p_a, p_b) + \min(1 - p_a, 1 - p_b) - p_e}{1 - p_e}.
\end{equation*}
For a white-box detector, the rates come from its learned threshold.

Let condition $c$ hold $n_c$ analysed rollouts, with rates $p_{a,c}$ and $p_{b,c}$. The
stratified $\kappa$ is $\sum_c n_c (p_{o,c} - p_{e,c}) / \sum_c n_c (1 - p_{e,c})$. Its
$\kappa_{\max}$ replaces $p_{o,c}$ with the largest attainable agreement in each condition. The
conditional-independence $\kappa$ is the overall $\kappa$ with $p_o$ replaced by
$\sum_c n_c [p_{a,c} p_{b,c} + (1 - p_{a,c})(1 - p_{b,c})]/N$, the agreement expected if $a$
and $b$ flagged independently within each condition. Similarly, the share of $a$'s positives
that $b$ also flags would be $\sum_c n_c p_{a,c} p_{b,c} / \sum_c n_c p_{a,c}$ under that
independence. The test-retest $\kappa$ is the stratified $\kappa$ over every ordered pair of
repetitions of one prompt. Apart from $d'$ and the overall $\kappa$, all these statistics are
post hoc.

\section{Related work}
\label{app:related}
Evaluation awareness appears in the reasoning of frontier models
\citep{apollo2025claude,schoen2025stress}, and anti-scheming stress tests give causal evidence
that it lowers covert behaviour \citep{schoen2025stress}. \citet{vishwarupe2026evaluation} prove that marginal
evaluation scores cannot identify the behavioural difference between recognised evaluation and
deployment. Studies of detection find that awareness is linearly decodable
\citep{chaudhary2025evaluation,heidari2026evaluation} and can be steered in open-weight reasoning models
\citep{abdelnabi2025hawthorne}. In open-weight models, however, it aligns only partly with
verbalisation \citep{heidari2026evaluation}, and detection couples only weakly with behaviour
and controllability \citep{nayan2026evaluation}.
Closest to our study, \citet{grover2026can} find that text baselines match a linear probe, and
conclude that a content-independent representation of evaluation awareness is yet to be shown.

\section{Additional results}
\label{app:results}
\begin{table}[htbp]
  \centering
  \footnotesize
  \setlength{\tabcolsep}{4pt}
  % Generated by scripts/paper/make_numbers.py from reports/full/analysis.json. Do not edit.
\begin{tabularx}{\linewidth}{@{}X ccc ccc@{}}
\toprule
 & \multicolumn{3}{c}{Gemma 4 31B} & \multicolumn{3}{c}{Qwen 3.6 27B} \\
\cmidrule(lr){2-4}\cmidrule(l){5-7}
Detector & $d'$ [95\% CI] & Hit & FA & $d'$ [95\% CI] & Hit & FA \\
\midrule
Self-report & 1.06 [0.88, 1.22] & 58.3 & 19.8 & 1.29 [1.11, 1.45] & 52.6 & 11.1 \\
Third-person & 1.17 [1.01, 1.32] & 80.0 & 37.1 & 1.60 [1.40, 1.79] & 74.7 & 17.6 \\
CoT judge & 1.17 [0.79, 1.65] & 5.2 & 0.3 & 0.68 [0.28, 1.36] & 3.3 & 0.6 \\
Activation probe & \multicolumn{3}{c}{excluded} & 1.24 [1.00, 1.48] & 62.1 & 17.6 \\
Logit monitor & \multicolumn{3}{c}{excluded} & 0.68 [0.53, 0.81] & 75.1 & 50.0 \\
\bottomrule
\end{tabularx}

  \caption{$d'$ per detector, all cued rollouts against baseline (cued: F1 to F7, in
  single-cue and multi-cue conditions). Hit and false-alarm (FA) rates in per cent, with the
  log-linear correction. Intervals on $d'$: cluster bootstrap over tasks. Gemma's white-box
  detectors are excluded (Appendix~\ref{app:gemma}).}
  \label{tab:sensitivity}
\end{table}

\subsection{Agreement over all analysed rollouts}
\label{app:overall}
Had self-report and third-person flagged independently within each cue condition, their overall
$\kappa$ would be \cv{gemma/sr-tp/kappaindep} on Gemma and \cv{qwen/sr-tp/kappaindep} on Qwen,
against the observed \cv{gemma/sr-tp/kappa} and \cv{qwen/sr-tp/kappa} (Table~\ref{tab:agreement}). In contrast, the Qwen
probe's overall $\kappa$ with each other detector is at most \cv{qwen/tp-probe/kappafrac}\% of
$\kappa_{\max}$, and it never exceeds the conditional-independence $\kappa$.

\begin{table}[htbp]
  \centering
  \footnotesize
  \begin{tabular*}{\linewidth}{@{\extracolsep{\fill}}l cccc@{}}
    \toprule
    Pair & $\kappa$ [95\% CI] & $\kappa_{\max}$ & \% of $\kappa_{\max}$ [95\% CI] & \shortstack{Cond.\ indep.\\$\kappa$} \\
    \midrule
    Gemma: self-report, third-person & \cv{gemma/sr-tp/kappa} [\cv{gemma/sr-tp/kappalo}, \cv{gemma/sr-tp/kappahi}] & \cv{gemma/sr-tp/kappamax} & \cv{gemma/sr-tp/kappafrac} [\cv{gemma/sr-tp/kappafraclo}, \cv{gemma/sr-tp/kappafrachi}] & \cv{gemma/sr-tp/kappaindep} \\
    Qwen: self-report, third-person & \cv{qwen/sr-tp/kappa} [\cv{qwen/sr-tp/kappalo}, \cv{qwen/sr-tp/kappahi}] & \cv{qwen/sr-tp/kappamax} & \cv{qwen/sr-tp/kappafrac} [\cv{qwen/sr-tp/kappafraclo}, \cv{qwen/sr-tp/kappafrachi}] & \cv{qwen/sr-tp/kappaindep} \\
    Qwen: self-report, logit monitor & \cv{qwen/sr-logit/kappa} [\cv{qwen/sr-logit/kappalo}, \cv{qwen/sr-logit/kappahi}] & \cv{qwen/sr-logit/kappamax} & \cv{qwen/sr-logit/kappafrac} [\cv{qwen/sr-logit/kappafraclo}, \cv{qwen/sr-logit/kappafrachi}] & \cv{qwen/sr-logit/kappaindep} \\
    Qwen: third-person, logit monitor & \cv{qwen/tp-logit/kappa} [\cv{qwen/tp-logit/kappalo}, \cv{qwen/tp-logit/kappahi}] & \cv{qwen/tp-logit/kappamax} & \cv{qwen/tp-logit/kappafrac} [\cv{qwen/tp-logit/kappafraclo}, \cv{qwen/tp-logit/kappafrachi}] & \cv{qwen/tp-logit/kappaindep} \\
    Qwen: third-person, probe & \cv{qwen/tp-probe/kappa} [\cv{qwen/tp-probe/kappalo}, \cv{qwen/tp-probe/kappahi}] & \cv{qwen/tp-probe/kappamax} & \cv{qwen/tp-probe/kappafrac} [\cv{qwen/tp-probe/kappafraclo}, \cv{qwen/tp-probe/kappafrachi}] & \cv{qwen/tp-probe/kappaindep} \\
    Qwen: self-report, probe & \cv{qwen/sr-probe/kappa} [\cv{qwen/sr-probe/kappalo}, \cv{qwen/sr-probe/kappahi}] & \cv{qwen/sr-probe/kappamax} & \cv{qwen/sr-probe/kappafrac} [\cv{qwen/sr-probe/kappafraclo}, \cv{qwen/sr-probe/kappafrachi}] & \cv{qwen/sr-probe/kappaindep} \\
    Qwen: probe, logit monitor & \cv{qwen/probe-logit/kappa} [\cv{qwen/probe-logit/kappalo}, \cv{qwen/probe-logit/kappahi}] & \cv{qwen/probe-logit/kappamax} & \cv{qwen/probe-logit/kappafrac} [\cv{qwen/probe-logit/kappafraclo}, \cv{qwen/probe-logit/kappafrachi}] & \cv{qwen/probe-logit/kappaindep} \\
    \bottomrule
  \end{tabular*}
  \caption{Overall agreement on the analysed rollouts. The overall $\kappa$ is the planned
  statistic; $\kappa_{\max}$, the percentage and Cond.\ indep.\ are post hoc. Intervals come
  from the cluster bootstrap over tasks. Cond.\ indep.: the overall $\kappa$ expected if the
  two detectors flagged independently within each cue condition.}
  \label{tab:agreement}
\end{table}

\subsection{Holding the task variant fixed}
\label{app:variant}
On baseline prompts, the safety variant raises the positive rates of self-report and
third-person (Table~\ref{tab:variant}). We therefore also stratified by cue condition and
variant together (post hoc). The pair then reaches \cv{gemma/sr-tp/kappacvfrac}\%
[\cv{gemma/sr-tp/kappacvfraclo}, \cv{gemma/sr-tp/kappacvfrachi}] of $\kappa_{\max}$ on Gemma and
\cv{qwen/sr-tp/kappacvfrac}\% [\cv{qwen/sr-tp/kappacvfraclo}, \cv{qwen/sr-tp/kappacvfrachi}] on
Qwen, and $\kappa_{\max}$ falls to \cv{gemma/sr-tp/kappacvmax} and \cv{qwen/sr-tp/kappacvmax}.
Unequal positive rates therefore still cause most of the disagreement. If the safety variant
itself signals an evaluation \citep{li2026decomposing}, these values understate the agreement.

\begin{table}[htbp]
  \centering
  \small
  \begin{tabular}{@{}l cc cc@{}}
    \toprule
    & \multicolumn{2}{c}{Gemma 4 31B} & \multicolumn{2}{c}{Qwen 3.6 27B} \\
    \cmidrule(lr){2-3}\cmidrule(l){4-5}
    Detector & Safety & Capability & Safety & Capability \\
    \midrule
    Self-report & \cv{gemma/sr/basesafety} & \cv{gemma/sr/basecap} & \cv{qwen/sr/basesafety} & \cv{qwen/sr/basecap} \\
    Third-person & \cv{gemma/tp/basesafety} & \cv{gemma/tp/basecap} & \cv{qwen/tp/basesafety} & \cv{qwen/tp/basecap} \\
    \bottomrule
  \end{tabular}
  \caption{Per cent of baseline rollouts flagged, by task variant (uncorrected rates).}
  \label{tab:variant}
\end{table}

\subsection{Shared positives}
\label{app:nesting}
Of the rollouts that self-report flags, third-person also flags \cv{gemma/sr-tp/nest}\% on
Gemma and \cv{qwen/sr-tp/nest}\% on Qwen. Under independence within cue conditions, these
shares would be \cv{gemma/sr-tp/nestindep}\% and \cv{qwen/sr-tp/nestindep}\%.
